% Width-one geometry and the separately specified anisotropic prior.
% Requires amsmath, amssymb, amsthm; lemma/proposition/corollary environments.
\section{Exact width-one geometry at the start of an AND layer}
\label{sec:and-width-one}

This section concerns a width-one network with one hidden tanh unit,
\[
 z_{a,b}(x)=b\tanh\!\left(a_0+\sum_{i=1}^d a_i x_i\right),
 \qquad x\in\{-1,1\}^d,
 \qquad
 \kappa(a,b)=d\sum_{i=0}^d a_i^2+b^2.
\]
The prior is the centered Gaussian with density proportional to
$\exp[-\kappa(a,b)/(2\sigma^2)]$, conditioned before training on
$|z_{a,b}(x)|>1/2$ at every cube input.
For a Boolean threshold function $g$, define its affine unit-margin cost
\[
 Q(g)=\min\left\{\sum_{i=0}^d a_i^2:
       g(x)\left(a_0+\sum_{i=1}^d a_i x_i\right)\ge1
       \text{ for every }x\right\}.
\]
The minimum exists: the feasible set is a nonempty closed polyhedron
and the objective is coercive. Non-threshold functions are not represented
by this width-one architecture.

\begin{lemma}[Reduction to affine margin cost]\label{lem:and-affine-reduction}
For every represented Boolean function $g$, its minimum strict-margin
neural energy is
\[
 E(g)=\Phi_d(Q(g)),\qquad
 \Phi_d(Q)=\min_{u>0}\left\{dQu^2+\frac14\coth^2u\right\}.
\]
The function $\Phi_d$ is strictly increasing on $Q>0$.
\end{lemma}
\begin{proof}
Replacing $(a,b)$ by $(-a,-b)$ preserves the real output and energy, so
we can take $b>0$. A strict output margin requires $b>1/2$. Put
$u=\operatorname{arctanh}(1/(2b))>0$. The constraints imply
$g(x)(a_0+\sum_i a_i x_i)>u$, so
$\sum_i a_i^2\ge u^2Q(g)$ and $b^2=\frac14\coth^2u$.
Conversely, multiply a unit-margin minimizer for $g$ by $(1+\epsilon)u$
and take $b=\frac12\coth u$. This gives strict output margin, with
energy tending to the displayed expression as $\epsilon\downarrow0$.
This proves equality of the infima. For $Q>0$ the objective tends to
infinity at both endpoints $u\downarrow0$ and $u\to\infty$, so its
minimum is attained at some $u_Q>0$. If $Q_2>Q_1$, then
\[
 \Phi_d(Q_2)
 =dQ_1u_{Q_2}^2+\tfrac14\coth^2u_{Q_2}
   +d(Q_2-Q_1)u_{Q_2}^2
 >\Phi_d(Q_1).
\]
\end{proof}

\paragraph{Small-variance function probabilities.}
For this finite-dimensional Gaussian law, the function event
$U_g=\{(a,b):g(x)z_{a,b}(x)>1/2\ \forall x\}$ satisfies
\begin{equation}\label{eq:affine-rate}
 \lim_{\sigma\downarrow0}\sigma^2\log\mathbb P_\sigma(U_g)
 =-\frac12 E(g).
\end{equation}
Indeed, the lower Gaussian large-deviation bound applies to the open
set $U_g$, and the upper bound applies to its closure. Their energy
infima coincide: any parameter vector with all signed outputs at
least $1/2$ can be approximated from within $U_g$ by multiplying the
readout by $1+\epsilon$. The standard finite-dimensional Gaussian
rate function is half the quadratic energy. The argument holds
unchanged for every positive definite diagonal quadratic energy.
There are only finitely many represented Boolean functions. Therefore,
after conditioning on any fixed label history, all strong-prior
posterior mass concentrates on its minimum-energy functions; if the
minimum is unique, its posterior probability tends to one.

Write $|x|_+=|\{i:x_i=+1\}|$. Fix $2\le k\le d$ and consider the
history that labels $1^d$ positive and every input with $|x|_+<k$
negative. These labels agree with the AND target, which is positive
only at $1^d$. No labels at other inputs are imposed yet.

\begin{lemma}[Essential affine coefficients]\label{lem:and-essential}
In a minimum-squared-norm unit-margin representation of a threshold
function, every nonzero input coefficient is essential to that function,
and every essential input coefficient has absolute value at least one.
\end{lemma}
\begin{proof}
If the function is independent of $x_i$, then, for each fixed assignment
of the other coordinates, both signs of $x_i$ have the same label.
Averaging their two signed margin inequalities shows that setting
$a_i=0$ preserves every unit-margin constraint. It strictly decreases
squared norm if $a_i\ne0$. If $x_i$ is essential, two neighboring cube
points differing only in $x_i$ have opposite labels. Their affine scores
have opposite signs and magnitudes at least one; their difference has
magnitude $2|a_i|$. Hence $|a_i|\ge1$.
\end{proof}

\begin{proposition}[All minimizers at a complete lower-layer prefix]
\label{prop:and-phase-minimizers}
Among all width-one Boolean functions consistent with this history,
the smallest affine unit-margin cost is
\[
 Q_{d,k}=
 \begin{cases}
  2k-1,&d\ge2k-1,\\
  d+(2k-1-d)^2,&d<2k-1.
 \end{cases}
\]
If $d\ge2k-1$, the minimizing functions are precisely
\[
 h_T(x)=+1
 \quad\Longleftrightarrow\quad
 |\{i\in T:x_i=+1\}|\ge k,
 \qquad |T|\in\{2k-1,2k-2\}.
\]
Their minimizing affine scores have coefficient one on $T$, zero
elsewhere, and bias zero for $|T|=2k-1$ or bias $-1$ for $|T|=2k-2$.
If $d<2k-1$, there is exactly one minimizing function: the symmetric
threshold $|x|_+\ge k$, with coefficients one and bias $d-2k+1$.
The corresponding minimum neural energy is $\Phi_d(Q_{d,k})$.
\end{proposition}
\begin{proof}
Choose a minimizing function and its minimum-norm unit-margin affine
representation. Lemma~\ref{lem:and-essential} applies. Let $p$ be the
number of positive input coefficients, write their sum as $P$, and let
$B$ be the sum of absolute values of the negative input coefficients.
In particular $P\ge p$.

One has $p\ge k$: the input which is positive exactly on the positive
coefficient coordinates has affine score at least that at $1^d$, hence
at least one. If $p<k$, the history requires this input to be negative,
a contradiction.

If $p\ge2k-1$, the positive coefficients alone contribute squared norm
at least $p\ge2k-1$. If $p\le2k-2$, average the negative margin
constraints over inputs that are positive on exactly $k-1$ of those
$p$ coordinates and negative everywhere else. This gives
\[
 a_0-P+B+\frac{2(k-1)}pP\le-1.
\]
Since $1-2(k-1)/p\le0$ and $P\ge p$,
\[
 a_0\le\left(1-\frac{2(k-1)}p\right)P-B-1
       \le p-2k+1-B\le p-2k+1<0.
\]
Consequently the total squared norm is at least
\[
 p+(2k-1-p)^2.
\tag{*}
\]
For integer $p\le2k-2$, this is at least $2k-1$, with equality only
at $p=2k-2$. Together with the preceding branch, equality at $2k-1$
forces either $p=2k-1$ with zero bias or $p=2k-2$ with bias $-1$.
It also forces every positive coefficient to equal one and excludes
all negative coefficients. The displayed functions realize these
bounds: all their affine scores are odd integers, and their positive
inputs contain at least $k$ positive coordinates.

If $d<2k-1$, every admissible $p$ lies between $k$ and $d\le2k-2$.
The expression in (*) strictly decreases as integer $p$ increases
in this range. It is therefore minimized only at $p=d$, giving
$d+(2k-1-d)^2$. Equality forces all coefficients to equal one and
the bias to equal $d-2k+1$. This score also takes odd integer values
and realizes the required symmetric threshold. The energy conclusion
follows from Lemma~\ref{lem:and-affine-reduction}.
\end{proof}

\paragraph{Why the circuit conjunction characterization fails here.}
A conjunction of $k$ positive literals has
\[
 Q(\operatorname{AND}_k)=k+(k-1)^2=k^2-k+1.
\]
To verify the lower bound, remove irrelevant coefficients, use
essentiality to make each remaining positive coefficient at least one,
and average the $k$ negative constraints obtained by flipping one of
the $k$ selected coordinates. This forces the bias to be at most
$-(k-1)$. Unit coefficients and that bias attain equality.
For $k\ge3$ and $d\ge2k-1$,
\[
 Q(\operatorname{AND}_k)-(2k-1)=(k-1)(k-2)>0.
\]
Thus every $k$-conjunction is strictly more expensive than the actual
minimizers at this prefix. The globally smallest circuit functions
at the analogous circuit prefix are $k$-conjunctions; the Gaussian
neural energy gives a different answer already in this one-unit model.

\paragraph{Exact masses in this particular minimizing class.}
Equality of minimum energy alone does not imply equal Gaussian function
probabilities. Here, however, an additional exact symmetry does imply it.
Ordinary input permutations make all functions with a fixed support
size equally probable. To connect the two support sizes, take
$j\in T$ with $|T|=2k-1$ and transform the affine parameters by
\[
 a'_0=-a_j,\qquad a'_j=a_0,\qquad a'_i=a_i\quad(i\ne j).
\]
This signed column swap preserves the Gaussian law and energy.
For $x\in\{-1,1\}^d$, set
\[
 y_j=-x_j,\qquad y_i=x_ix_j\quad(i\ne j).
\]
Then $a'_0+\sum_i a'_ix_i=x_j(a_0+\sum_i a_iy_i)$, so oddness of
tanh gives $z'(x)=x_jz(y)$. The map $x\mapsto y$ is a bijection,
and hence the global-margin event is preserved. Applied to majority
on $T$, the transformed sign is
\[
 x_j\operatorname{sign}\!\left(\sum_{i\in T}y_i\right)
 =\operatorname{sign}\!\left(-1+\sum_{i\in T\setminus\{j\}}x_i\right),
\]
which is precisely the threshold of $k$ positives on the even support
$T\setminus\{j\}$. Thus all functions in the minimizing class have
exactly equal prior mass at every positive variance. Equal mass persists
among whichever of them survive later label conditioning.

By the fixed-width Gaussian rate theorem and finiteness of the Boolean
function class, the strong-prior posterior at this prefix converges to
uniform sampling from these minimizers. At positive variance it also
assigns mass to higher-energy functions; it is not exactly this uniform
minimum-class rule. In other histories, equal energies need not be
accompanied by the symmetry just proved.

\begin{corollary}[The first minimum-energy tier cannot itself cause an
exponential phase]\label{cor:and-first-tier}
Suppose $d\ge2k-1$ and start from the complete lower-layer prefix above.
Let
\[
 M_{d,k}=\binom d{2k-1}+\binom d{2k-2}.
\]
For the width-one strong-prior limiting posterior, the expected number
of sampled mistakes before this minimum-energy tier is exhausted is at
most $1+\log M_{d,k}\le1+(d+1)\log2$. The bound includes the round which exhausts
it. It holds under every adaptive nonanticipating continuation. If
$d<2k-1$ and $k<d$, its sole initial minimum is eliminated by the first
weight-$k$ negative example, causing one mistake on that example. If
$k=d$, the prefix has already revealed the entire cube and no further
mistakes are possible.
\end{corollary}
\begin{proof}
While at least one of the original minimizers survives, it still has
energy $\Phi_d(Q_{d,k})$, so every strong-prior posterior is uniform
on the surviving original minimizers. Eliminating a fraction $q<1$
of a uniform finite class reduces its logarithmic size by at least $q$.
The final exhaustion costs at most one. Summing these conditional error
probabilities and taking expectations proves the claim. The singleton
case follows from the unique symmetric threshold's positive prediction
on every weight-$k$ input.
\end{proof}

\paragraph{Limits of the conclusion.}
The proposition characterizes the first energy tier at a complete
lower-layer prefix. After that tier is exhausted, higher-energy
functions can become optimal; their complete sequence is not
characterized here. Thus the corollary is not an upper bound for the
whole Hamming layer or for the entire AND-learning process.
All statements concern width one, one hidden tanh unit, and global
margin conditioning. For a growing AND target, width one need not
satisfy the separate theorem convention $N\ge C_L(f)$, because that
complexity also bounds the magnitude of every teacher weight. No
extension to all admissible widths, and no exponential lower bound
for the original Gaussian neural posterior, is asserted by this width-one calculation.

\subsection{An arbitrarily small anisotropy changes the width-one answer}
\label{sec:and-anisotropy}

The preceding equality of masses uses exact symmetry. The next result
changes the prior and is not a lower bound for the canonical isotropic
first-layer law. It also retains width one; it does not satisfy the
separate $N\ge C_L(f)$ convention for the growing AND target in general.
It does show that a small independent-Gaussian perturbation can turn the
same concrete neural minimizing class into an exponential mistake path.

\begin{proposition}[Near-isotropic independent-Gaussian variant]
\label{prop:and-anisotropic}
Let $k\ge2$ and $d\ge8k$. For every $\eta>0$, there are coefficients
$1\le\lambda_i\le1+\eta$, $0\le i\le d$, with the following property.
Use a width-one prior with independent centered Gaussian weights
\[
 \operatorname{Var}(a_i)=\frac{\sigma^2}{d\lambda_i},\qquad
 \operatorname{Var}(b)=\sigma^2,
\]
and the same all-input output-margin conditioning. Its quadratic energy is
\[
 \kappa_\lambda(a,b)=d\sum_{i=0}^d\lambda_i a_i^2+b^2.
\]
For the AND target, there is a fixed distinct-input order on
which the strong-prior posterior makes at least
\[
 \frac{\binom d k}{\binom{2k-1}k}\ge2^k
\]
posterior-sampling mistakes with probability one. Posterior majority is
wrong on the same rounds.

For every fixed $0<\zeta<1/2$, one sufficiently small positive variance
makes the actual Gaussian posterior's wrong-label probability at least
$1-\zeta$ on every one of these rounds. Its expected sampled mistakes
on them are therefore at least $(1-\zeta)2^k$, and its majority is
wrong on every one. In particular, $k=\lfloor d/8\rfloor$ gives
an exponential lower bound for this perturbed width-one model.
\end{proposition}
\begin{proof}
Start with the history of Proposition~\ref{prop:and-phase-minimizers}.
Let $\mathcal C$ be its minimizing class and $Q_0=2k-1$. Every other
consistent threshold function has strictly larger isotropic affine cost.
There are finitely many threshold functions, so their costs exceed $Q_0$
by a positive common gap $\Delta$.

For $\epsilon>0$, put $\lambda_i=1+\epsilon v_i$, where the numbers
$v_i\in(0,1]$ have distinct sums over distinct subsets of
$\{0,\ldots,d\}$; for example $v_i=2^{-i}$ has this property.
For each fixed threshold function $h$, let
\[
 Q_h(\epsilon)=\min\left\{\sum_i(1+\epsilon v_i)a_i^2:
             h(x)(a_0+\textstyle\sum_{i\ge1}a_ix_i)\ge1
             \text{ for all }x\right\}.
\]
Its isotropic optimizer $a_h$ is unique by strict convexity. The
optimizers at $\epsilon\downarrow0$ converge to $a_h$: boundedness
follows from comparison with $a_h$, and every limit point is an isotropic
optimizer. Comparing the perturbed and unperturbed objectives in both
directions consequently gives
\[
 Q_h(\epsilon)=Q_h(0)+\epsilon\sum_i v_i(a_h)_i^2+o(\epsilon).
\]
For $h\in\mathcal C$, the nonzero squared coordinates of $a_h$ are
exactly the indicator of its odd support $T$, or of its even support
$T$ together with the bias index $0$. These augmented supports are
pairwise different. Their derivative sums are therefore pairwise
distinct. Since $\mathcal C$ is finite, all its perturbed costs are
pairwise distinct for sufficiently small positive $\epsilon$.

For every $h\in\mathcal C$,
$Q_h(\epsilon)\le(1+\epsilon)Q_0$. All other initially consistent
functions have $Q_h(\epsilon)\ge Q_0+\Delta$, because $v_i>0$.
Taking also $\epsilon<\Delta/Q_0$ and $\epsilon\le\eta$ ensures that
every member of $\mathcal C$ remains cheaper than every outsider.
The affine-to-tanh reduction is unchanged, with $Q_h(\epsilon)$ in
place of $Q(h)$, so the same ordering holds for neural energies.

Until $\mathcal C$ is exhausted, the currently consistent function of
minimum energy is therefore a unique member of $\mathcal C$. Choose
any $k$-element subset of its support and query the input positive
exactly on that subset. The selected function predicts positive; the
AND target is negative because $k<d$. The query is new, since a
consistent selected function cannot be positive on a previously
revealed negative query. This step eliminates exactly those candidate
supports containing the queried $k$-set.

There are $\binom d{2k-1}$ candidates with odd support size $2k-1$.
A single $k$-set query eliminates at most
$\binom{d-k}{k-1}$ of these. Exhausting $\mathcal C$ therefore
requires at least
\[
 \frac{\binom d{2k-1}}{\binom{d-k}{k-1}}
 =\frac{\binom d k}{\binom{2k-1}k}
 \ge\frac{(d/k)^k}{4^k}\ge2^k
\]
queries, each a certain mistake in the strong-prior limit.
Use a fixed convention for choosing each $k$-set; the complete query
sequence can then be computed before the run and completed to a
permutation of the cube. It does not inspect any current random draw.

Finally, at each of these finitely many histories the minimum-energy
function is unique with a positive energy gap. The Gaussian rate
argument applies to every positive definite diagonal quadratic energy.
Thus, as the common variance tends to zero, the posterior mass of that
function tends to one. Taking the minimum of finitely many positive
variance thresholds makes its wrong-label probability at least
$1-\zeta$ at every selected round simultaneously. The input sequence
and labels are fixed, so posterior randomness at earlier rounds does
not alter this argument.
\end{proof}

\begin{corollary}[An admissible width-one majority target]
\label{cor:majority-anisotropic}
Let $d\ge17$ be odd and put $k=\lfloor d/8\rfloor$. The conclusion
of Proposition~\ref{prop:and-anisotropic} also holds for the target
\[
 f_{\rm maj}(x)=\operatorname{sign}\!\left(\sum_{i=1}^d x_i\right).
\]
This target has bounded-weight teacher complexity $C_2(f_{\rm maj})=1$,
so the learner width $N=1$ satisfies $N\ge C_2(f_{\rm maj})$.
For every $\eta>0$, an independent-Gaussian first-layer covariance
perturbation with $1\le\lambda_i\le1+\eta$ therefore admits a fixed
input permutation with at least $2^{\lfloor d/8\rfloor}$ certain
strong-prior mistakes. At all sufficiently small positive variances,
the expected sampled mistakes on those rounds are at least
$(1-\zeta)2^{\lfloor d/8\rfloor}$ for any prescribed
$0<\zeta<1/2$, and posterior majority is wrong on all of them.
\end{corollary}
\begin{proof}
The network $z(x)=\tanh(\sum_i x_i)$ has one hidden unit, zero bias,
and all input and readout weights equal to one. Since $d$ is odd,
$|\sum_i x_i|\ge1$ on the cube, so its signed output margin is at
least $\tanh(1)>1/2$. Thus its bounded-weight teacher complexity is
one.

The construction in Proposition~\ref{prop:and-anisotropic} uses only
the positive label at $1^d$ and negative labels at inputs of Hamming
weight at most $k$. Because $k<d/2$, all these labels agree with
$f_{\rm maj}$. Consequently every posterior and every forced query
in that construction is unchanged when the target is $f_{\rm maj}$.
After that finite sequence, reveal any remaining inputs with their
majority labels. The same Gaussian prior, energy gaps, and positive
variance thresholds prove the stated bounds.
\end{proof}

\paragraph{Scope of the perturbed examples.}
The parameters of the width-one target
$\tanh(\sum_i x_i-d+1)$ have magnitude $O(d)$ and strict output margin
$\tanh(1)>1/2$. Thus the example has a simple, polynomially bounded
neural teacher, but this AND realization does not establish the
bounded-weight width-admissibility convention.
Corollary~\ref{cor:majority-anisotropic} removes that width issue for
a different simple target. Both examples change the exact coordinate
symmetry of the first layer. Neither is an exponential lower bound
for the canonical isotropic prior; no extension of the AND lower
bound to all admissible larger learner widths is established.

\paragraph{High probability and ordinary iid learning.}
The finite-variance lower bound can also force all the designated
mistakes with probability at least $1-\delta$: if there are $T$ forced
rounds, apply the proposition with $\zeta<\delta/T$ and use a union
bound. Choosing also $\zeta<1/3$ makes every designated round a failure
of the strict $2/3$ confidence requirement.

The same model retains polynomial iid sample complexity. Its supported
Boolean functions are affine threshold functions. The arrangement of
$m=2^d$ homogeneous hyperplanes in the $(d+1)$-dimensional affine
coefficient space gives at most
$2\sum_{j=0}^d\binom{m-1}{j}\le2(d+1)m^d$
possible strict sign patterns. For iid examples from any fixed input
distribution, a union bound therefore implies that every consistent
function has error at most $\varepsilon$, except with probability
$\delta$, once
\[
 n\ge\frac{\log\bigl(2(d+1)2^{d^2}\bigr)+\log(1/\delta)}
              {\varepsilon}.
\]
This controls posterior-sampling test error for every prior on this
class. The test-distribution probability of inputs failing strict $2/3$ confidence is at
most three times the posterior-averaged error. Thus the anisotropic
example separates polynomial iid learning from exponential online
mistakes under an adversarially chosen fixed order; it does not merely
use a model that was already statistically hard under iid data.
